%% ma: L11-B16-0070
%% lop: 11
%% chuong: 5
%% bai: 16
%% dang: 11.16.3
%% loai: TLN
%% muc: VD
%% dap_an: "7"
%% nguon: "tai-lieu-hoc-tap-mon-toan-lop-11-theo-chuong-trinh-giao-duc-pho-thong-2018.pdf (Lớp Toán Thầy Duy, Phạm Lê Duy, THPT Chu Văn An), Chương 3 Bài 2, Phần 3, Câu 8 (tr.381)"
%% kien_thuc: "Lớp 11 Bài 16 (giới hạn của hàm số)"
%% trang_thai: cho-duyet
%% ly_do: "Dạng mới đề xuất (11.16.x), chờ giáo viên duyệt dạng; tài liệu gốc không có lời giải mẫu, đáp án do kho tự giải và kiểm bằng sympy."
\begin{ex}
Biến đổi giới hạn $\lim\limits_{x\to+\infty}\dfrac{(6-2x^4)^4\cdot(2-x^2)}{(x^2-3)^5\cdot(6-4x)}$ ta thu được kết quả dạng \[\lim\limits_{x\to+\infty}x^a\cdot\dfrac{\left(\dfrac6{x^4}-2\right)^4\cdot\left(\dfrac2{x^2}-1\right)}{\left(1-\dfrac3{x^2}\right)^5\cdot\left(\dfrac6x-4\right)}\] với $a$ là số tự nhiên. Xác định giá trị của $a$.
\shortans{7}
\loigiai{\sloppy Tử có bậc $4\cdot4+2=18$, mẫu có bậc $2\cdot5+1=11$. Đưa $x^{16}$, $x^2$ ở tử và $x^{10}$, $x$ ở mẫu ra ngoài: $x^{18-11}=x^7$. Vậy $a=7$.}
\end{ex}
